% Эллипс и вспомогательная окружность: a = 18 мм, e = 0.8, E = 120 градусов.
% Q и КА имеют одинаковую абсциссу; F — правый фокус, П — перицентр.
\begin{tikzpicture}[course diagram,x=1mm,y=1mm]
  \path[use as bounding box] (-29,-19) rectangle (25,21);
  \coordinate (C) at (0,0);
  \coordinate (F) at (14.4,0);
  \coordinate (P) at (18,0);
  \coordinate (Q) at (-9,15.58846);
  \coordinate (K) at (-9,9.35307);

  % Вспомогательная окружность сохраняет истинные пропорции.
  \draw[coursemuted,line width=.5pt,dash pattern=on 2pt off 2pt]
    (C) circle[radius=18];
  \draw[orbit line] (C) ellipse[x radius=18,y radius=10.8];
  \draw[aux line] (-19.5,0)--(21.5,0);
  \draw[coursemuted,line width=.6pt] (C)--(Q);
  \draw[coursemuted,line width=.5pt,dashed] (-9,0)--(Q);
  \draw[vector line] (F)--(K);

  % E отсчитывается при центре, nu — при фокусе от одного направления оси.
  \draw[angle line] (4.5,0)
    arc[start angle=0,end angle=120,radius=4.5];
  \node[text=courseaccent,inner sep=.4mm,fill=coursepaper]
    at (1.4,6.9) {$E$};
  \draw[angle line] (20.4,0)
    arc[start angle=0,end angle=158.21492,radius=6];
  \node[text=courseaccent,inner sep=.4mm,fill=coursepaper]
    at (15.8,8.1) {$\nu$};

  \fill[courseink] (C) circle[radius=.55];
  \fill[courseink] (F) circle[radius=.65];
  \fill[courseaccent] (P) circle[radius=.6];
  \draw[courseink,fill=coursepaper,line width=.65pt]
    (Q) circle[radius=.7];
  \fill[courseaccent] (K) circle[radius=.8];
  \node[anchor=north,inner sep=1mm] at (C) {$C$};
  \node[anchor=north,inner sep=1mm] at (F) {$F$};
  \node[anchor=north west,inner sep=1mm] at (P) {П};
  \node[anchor=south east,inner sep=1mm] at (Q) {$Q$};
  \draw[coursemuted,line width=.5pt] (-20,9.35307)--(-10.1,9.35307);
  \node[anchor=east,inner sep=1mm] at (-20,9.35307) {КА};
\end{tikzpicture}
